\documentclass[11pt,letterpaper]{article}
\usepackage[utf8]{inputenc}
\usepackage[french]{babel}
\usepackage{amsmath,amssymb,amsthm}
\usepackage{geometry}
\geometry{margin=2.5cm}
\usepackage{enumitem}
\usepackage{hyperref}
\usepackage{booktabs}
\usepackage[
	backend=biber,
	style=authoryear,
	natbib=true,
	sorting=nyt
]{biblatex}
\IfFileExists{STT-4300-references.bib}{%
	\addbibresource{STT-4300-references.bib}%
}{%
	\addbibresource{STT-4300/STT-4300-references.bib}%
}

\providecommand{\citeN}{\citet}
\providecommand{\citeNP}{\citealt}

\newtheorem{definition}{Définition}[section]
\newtheorem{propriete}{Propriété}[section]
\newtheorem{remarque}{Remarque}[section]
\newtheorem{exemple}{Exemple}[section]

\title{\textbf{Notes de cours --- STT-4300}\\Modèles linéaires généralisés et additifs}
\author{}
\date{}

\begin{document}
\maketitle
\tableofcontents
\newpage

%% ============================================================
%% CHAPITRE 1 : Introduction et motivation
%% ============================================================
\section{Introduction et motivation}

\subsection{Formulation générale}

Tous les modèles étudiés dans ce cours partagent la même formulation générale~:
$$g(E(Y_i)) = g(\mu_i) = \eta_i = \begin{cases}
\beta_0 + \beta_1 x_{i1} + \cdots + \beta_p x_{ip} & \text{(linéaire)}\\
\alpha + f_1(x_{i1}) + \cdots + f_p(x_{ip}) & \text{(additif)}
\end{cases}$$
où $g$ est la \textbf{fonction de lien} (\textit{link function}) et $\eta_i$ est le \textbf{prédicteur linéaire} (ou additif).

\subsection{Notation}

\textbf{Variable réponse~:} observations $y_1,\ldots,y_n$ pour $n$ individus, correspondant aux variables aléatoires $Y_1,\ldots,Y_n$.

\textbf{Propriétés distributionnelles~:} $Y_i \sim F_{\theta_i}$, avec $E(Y_i) = \mu_i$ et $Var(Y_i) = v_i = v(\mu_i)$.

\textbf{Variables explicatives~:} $x_i = (x_{i0}, x_{i1},\ldots,x_{ip})^T$, où $x_{i0}=1$ (ordonnée à l'origine). La matrice de plan (\textit{design matrix}) est
$$X = \begin{pmatrix} x_1^T \\ \vdots \\ x_n^T \end{pmatrix} = \begin{pmatrix} x_{10} & x_{11} & \cdots & x_{1p} \\ \vdots & \vdots & & \vdots \\ x_{n0} & x_{n1} & \cdots & x_{np}\end{pmatrix}.$$

\subsection{Principes de la modélisation statistique}

\begin{enumerate}
\item \textbf{Analyse exploratoire des données~:} vérifier la qualité des données et aider à la formulation du modèle (échelle de mesure, forme de la distribution, associations entre variables).
\item \textbf{Formulation du modèle~:} distribution de probabilité de $Y$ et lien avec les variables explicatives.
\item \textbf{Estimation des paramètres~:} maximum de vraisemblance, moindres carrés, méthode des moments, etc.
\item \textbf{Résidus et vérification du modèle~:} inspection graphique des résidus, tests d'adéquation.
\item \textbf{Inférence et interprétation~:} tester les hypothèses pour obtenir un modèle parcimonieux, puis interpréter les paramètres.
\end{enumerate}

\newpage
%% ============================================================
%% CHAPITRE 2 : Modèles linéaires généralisés (MLG)
%% ============================================================
\section{Modèles linéaires généralisés (MLG)}

\subsection{Définition}

Le modèle linéaire classique $Y_i = x_i^T \beta + \epsilon_i$ avec $\epsilon_i \sim \mathcal{N}(0,\sigma^2)$ est un cas particulier. Les MLG généralisent en permettant~:
\begin{enumerate}
\item des distributions de la variable réponse autres que la loi normale (même discrètes ou catégorielles),
\item une relation entre la réponse et les variables explicatives qui n'est pas nécessairement linéaire simple.
\end{enumerate}

La forme générale d'un MLG est~:
$$\boxed{g(E(Y_i)) = g(\mu_i) = \eta_i = x_i^T \beta}$$
où $g$ est la fonction de lien et $Var(Y_i) = v(\mu_i)$.

\subsection{Famille exponentielle}

\begin{definition}
Une densité (ou fonction de probabilité) appartient à la \textbf{famille exponentielle} si elle s'écrit~:
$$f(y;\theta,\phi) = \exp\left(A\,\frac{y\theta - b(\theta)}{\phi} + c\Big(y,\frac{\phi}{A}\Big)\right),$$
pour des fonctions spécifiques $b(\cdot)$ et $c(\cdot)$ et une constante $A$. Si $\phi$ est connu, c'est une famille exponentielle avec \textbf{paramètre canonique} $\theta$.
\end{definition}

Une fonction de lien $g$ est dite \textbf{canonique} si $\theta(\mu) = \eta$, c'est-à-dire $g(\cdot) = (b')^{-1}(\cdot)$.

\begin{propriete}
Pour la famille exponentielle~:
\begin{align}
E(Y) &= \mu = b'(\theta) \label{eq:esp}\\
Var(Y) &= \frac{b''(\theta)\,\phi}{A} \label{eq:var}
\end{align}
\end{propriete}

\textbf{Démonstration.} On utilise les identités classiques~:
$$E\left(\frac{\partial l}{\partial \theta}\right) = 0 \quad \text{et} \quad E\left(\frac{\partial^2 l}{\partial \theta^2}\right) + E\left(\left(\frac{\partial l}{\partial \theta}\right)^2\right) = 0.$$
Comme $\frac{\partial l}{\partial \theta} = A\frac{y - b'(\theta)}{\phi}$ et $\frac{\partial^2 l}{\partial \theta^2} = -A\frac{b''(\theta)}{\phi}$, on obtient les résultats \eqref{eq:esp} et \eqref{eq:var}.

\bigskip

\begin{center}
\begin{tabular}{|l|c|c|c|c|}
\hline
Distribution & $A$ & $\phi$ & $b(\theta)$ & Lien canonique \\
\hline
$\mathcal{N}(\mu,\sigma^2)$ & $1$ & $\sigma^2$ & $\theta^2/2$ & Identité \\
$\mathcal{P}(\lambda)$ & $1$ & $1$ & $\exp(\theta)$ & $\log$ \\
$\mathcal{B}(m,p)/m$ & & & & Logit \\
$\Gamma(\mu,\nu)$ & & & & $1/\mu$ \\
\hline
\end{tabular}
\end{center}

\subsection{Composantes d'un MLG}

Un MLG comporte \textbf{trois composantes}~:
\begin{enumerate}
\item Des réponses indépendantes $Y_1,\ldots,Y_n$ de la même distribution de la famille exponentielle.
\item Un ensemble de paramètres $\beta$ et de variables explicatives $x_i^T = (1, x_{i1},\ldots,x_{ip})$.
\item Une fonction de lien $g$ monotone et dérivable telle que $g(\mu_i) = x_i^T \beta = \eta_i$.
\end{enumerate}

La transformation via la fonction de lien assure que $\mu_i$ estimé se situe dans l'espace admissible des valeurs (par ex.\ $(0,1)$ pour Bernoulli, $(0,\infty)$ pour Poisson).

\subsection{Exemples de MLG}

\subsubsection{Loi normale~: $Y_i \sim \mathcal{N}(\mu_i, \sigma^2)$}
Lien canonique~: identité $g(\mu_i) = \mu_i$. Utilisée pour les données continues symétriques.

\subsubsection{Loi binomiale~: $Y_i \sim \mathcal{B}(m_i, p_i)$}
Trois fonctions de lien courantes (restreignent $\mu_i \in (0,1)$)~:
\begin{itemize}
\item \textbf{Logit} (canonique)~: $g(p_i) = \log\left(\frac{p_i}{1-p_i}\right)$
\item \textbf{Probit}~: $g(p_i) = \Phi^{-1}(p_i)$
\item \textbf{Log-log complémentaire}~: $g(p_i) = \log(-\log(1-p_i))$
\end{itemize}

Les valeurs $o_i = p_i/(1-p_i)$ sont les \textbf{cotes} (\textit{odds}). Le lien logit permet l'interprétation sur l'échelle des cotes~: un changement unitaire de $x_j$ multiplie les cotes par $\exp(\beta_j)$.

\subsubsection{Loi de Poisson~: $Y_i \sim \mathcal{P}(\lambda_i)$}
Lien canonique~: $g(\lambda_i) = \log(\lambda_i)$. Utilisée pour les données de dénombrement.

\subsubsection{Loi gamma~: $Y_i \sim \Gamma(\mu_i, \nu)$}
Utilisée pour les données continues asymétriques (revenus, coûts). Le lien courant est $g(\mu_i) = \log(\mu_i)$. Le lien canonique $1/\mu_i$ est plus rarement utilisé.

\subsection{Extensions~: loi multinomiale}

\subsubsection{Données nominales}
On lie les prédicteurs aux probabilités par~:
$$\log\left(\frac{p_{ik}}{p_{i1}}\right) = x_i^T \beta_k, \quad k=2,\ldots,K.$$

\subsubsection{Données ordinales}
Modèle à \textbf{cotes proportionnelles}~:
$$\log\left(\frac{\gamma_k(x_i)}{1-\gamma_k(x_i)}\right) = \alpha_k - x_i^T \beta, \quad k=1,\ldots,K-1,$$
où $\gamma_k(x_i) = P(Y_i \leq k \mid x_i)$. Le rapport de cotes $\frac{\text{odds}(Y \leq k \mid x)}{\text{odds}(Y \leq k \mid \tilde{x})} = \exp(-(x-\tilde{x})^T\beta)$ est indépendant de $\alpha_k$.

\subsection{Estimation des paramètres~: maximum de vraisemblance}

La log-vraisemblance de la famille exponentielle est~:
$$l(\theta_1,\ldots,\theta_n; y) = \frac{A}{\phi}\sum_{i=1}^n y_i\theta_i - \frac{A}{\phi}\sum_{i=1}^n b(\theta_i) + \sum_{i=1}^n c\Big(y_i, \frac{\phi}{A}\Big).$$

On utilise la relation $g(\mu_i) = x_i^T\beta$ et on maximise par rapport à $\beta$ en résolvant les équations de score. Aucune forme fermée en général.

\subsubsection{Newton-Raphson et méthode de scoring}
Pour maximiser $l(\beta;y)$~:
\begin{itemize}
\item \textbf{Newton-Raphson}~: $\beta^{(m)} = \beta^{(m-1)} - [U'^{(m-1)}]^{-1} U^{(m-1)}$
\item \textbf{Méthode de scoring}~: $\beta^{(m)} = \beta^{(m-1)} + [\mathcal{J}^{(m-1)}]^{-1} U^{(m-1)}$
\end{itemize}
où $U = \partial l/\partial \beta$ et $\mathcal{J} = -E(U') $ (information de Fisher).

\subsection{Algorithme MCPI (IWLS)}

L'algorithme des \textbf{moindres carrés pondérés itératifs} résout itérativement~:
$$X^T W^{(m-1)} X \, \beta^{(m)} = X^T W^{(m-1)} z^{(m-1)},$$
où $W$ est diagonale avec $w_{ii} = \frac{1}{Var(Y_i)}\left(\frac{\partial \mu_i}{\partial \eta_i}\right)^2$ et 
$$z_i^{(m-1)} = x_i^T \beta^{(m-1)} + (y_i - \mu_i^{(m-1)})\left(\frac{\partial \eta_i}{\partial \mu_i}\right)^{(m-1)}.$$

\begin{remarque}
Lorsque le lien canonique est utilisé, les équations de score se simplifient~: $U_j = \sum_i (y_i - \mu_i) x_{ij}$.
\end{remarque}

\subsection{Vérification du modèle}

\subsubsection{Déviance}

La \textbf{déviance réduite} est~:
$$D^*(y,\mu) = 2[l(y;y) - l(\hat{\beta};y)] = \sum_{i=1}^n d_i,$$
et la \textbf{déviance}~: $D(y,\mu) = \phi\, D^*(y,\mu)$.

De grandes valeurs de $D$ indiquent un mauvais ajustement. Sous certaines conditions, $D^* \sim \chi^2_{n-(p+1)}$ approximativement.

\textbf{Attention}~: $D$ est inutilisable quand $Y_i \sim \text{Bernoulli}(p_i)$.

\subsubsection{Statistique $X^2$ de Pearson}
$$X^2 = \sum_{i=1}^n \frac{(y_i - \hat{\mu}_i)^2}{v(\hat{\mu}_i)}.$$

\subsubsection{Mesures globales}
\begin{itemize}
\item Proportion de déviance expliquée~: $(D_{null} - D)/D_{null}$.
\item Extension du $R^2$ de Nagelkerke~: $\widetilde{R}^2 = \frac{1 - (L_{null}/L)^{2/n}}{1 - L_{null}^{2/n}}$.
\end{itemize}

\subsubsection{Résidus}
\begin{itemize}
\item \textbf{Résidus de Pearson}~: $r_{iP} = \frac{y_i - \hat{\mu}_i}{\sqrt{v(\hat{\mu}_i)}}$
\item \textbf{Résidus de déviance}~: $r_{iD} = \text{sign}(y_i - \hat{\mu}_i)\sqrt{d_i}$
\item \textbf{Résidus de déviance standardisés}~: $r_{iDS} = r_{iD}/\sqrt{1-h_{ii}}$
\item \textbf{Résidus quantiles randomisés}~: si $F$ continue, $r_{iQ} = \Phi^{-1}(F(y_i; \hat{\mu}_i, \hat{\phi}))$. Distribution exactement $\mathcal{N}(0,1)$ sous le vrai modèle.
\end{itemize}

\textbf{Note}~: $\sum r_{iP}^2 = X^2$ et $\sum r_{iD}^2 = D$.

Les résidus sont représentés~: vs valeurs ajustées, vs quantiles normaux (QQ), dans l'ordre des observations, vs chaque covariable continue.

\subsection{Intervalles de confiance et tests z}

De la normalité asymptotique $(\hat{\beta} - \beta) \sim \mathcal{N}(0, \tilde{\mathcal{J}}^{-1})$~:
$$IC_{1-\alpha}(\beta_i) = \left[\hat{\beta}_i \pm z_{1-\alpha/2} \sqrt{(\tilde{\mathcal{J}}^{-1})_{ii}}\right].$$

Les \textbf{intervalles de vraisemblance profilée} ne supposent pas la normalité de l'estimateur et fonctionnent mieux en petits échantillons (obtenus en inversant les tests du rapport de vraisemblance).

\textbf{Statistique z (ou t)}~: $z = \hat{\beta}_i / se(\hat{\beta}_i) \sim \mathcal{N}(0,1)$ sous $H_0: \beta_i = 0$. Si $\phi$ est inconnu, on utilise $t_{n-(p+1)}$.

\subsection{Tests d'hypothèses}

Pour tester $H_0: \beta = \tilde{\beta}$~:
\begin{description}
\item[Wald~:] $(\hat{\beta}-\tilde{\beta})^T \mathcal{J}(\tilde{\beta})(\hat{\beta}-\tilde{\beta}) \sim \chi^2_{p+1}$
\item[Score~:] $U(\tilde{\beta})^T \mathcal{J}(\tilde{\beta})^{-1} U(\tilde{\beta}) \sim \chi^2_{p+1}$
\item[Rapport de vraisemblance~:] $2[l(\hat{\beta};y) - l(\tilde{\beta};y)] \sim \chi^2_{p+1}$
\end{description}

\subsection{Sélection de variables}

\subsubsection{Différence de déviances}

Pour comparer un modèle $\mathcal{M}_{p+1}$ à un sous-modèle $\mathcal{M}_{p+1-q}$~:
$$\Delta D^* = D^*(y,\hat{\mu}^{p+1-q}) - D^*(y,\hat{\mu}^{p+1}) \sim \chi^2_q \quad \text{sous } H_0.$$

Si $\phi$ est inconnu~:
$$\frac{(D(y,\hat{\mu}^{p+1-q}) - D(y,\hat{\mu}^{p+1}))/q}{\hat{\phi}} \sim F_{q,\, n-(p+1)}.$$

\subsubsection{Critère d'Akaike (AIC)}
$$AIC = -2\,l(\hat{\beta}^d;y) + 2d.$$
Les modèles avec un AIC faible sont préférés. On peut utiliser l'AIC pas-à-pas lorsque $2^p$ est trop grand.

\subsection{Ce qui peut mal tourner}

\subsubsection{Surdispersion}
Si $Var(Y_i) > E(Y_i)$ (Poisson) ou $Var(Y_i) > m_i p_i(1-p_i)$ (binomiale), on parle de \textbf{surdispersion}. On peut~:
\begin{itemize}
\item Utiliser un modèle quasi-vraisemblance avec $Var(Y_i) = \sigma^2 v(\mu_i)$ ($\sigma^2$ estimé).
\item Utiliser la \textbf{loi binomiale négative}~: $Var(Y_i) = \mu_i(1 + \mu_i/\vartheta)$.
\item Utiliser la \textbf{loi bêta-binomiale} pour les données binomiales surdispersées.
\end{itemize}

\textbf{Modèle de taux}~: le terme $\log(\text{Aggregate})$ (offset) contrôle le nombre d'occasions (\textit{effort}).

\subsubsection{Séparation (quasi-)complète}
En régression binaire, si les réponses 0 et 1 peuvent être (presque) séparées par les covariables, l'EMV n'existe pas ou diverge. Corrections possibles~: régression logistique exacte, approche robuste, méthode de réduction du biais de Firth.

\subsubsection{Phénomène de Hauck-Donner}
La statistique $z$ de Wald peut être trompeuse quand les probabilités ajustées sont proches de 0 ou 1~: une petite valeur de $z$ peut correspondre à un $|\hat{\beta}|$ très grand avec une mauvaise approximation de Wald.

\newpage
%% ============================================================
%% CHAPITRE 3 : Excès de zéros
%% ============================================================
\section{Modélisation de l'excès de zéros dans les données de comptage}

Dans de nombreuses applications, la variable réponse contient plus de zéros que ce qu'impliquerait une loi de Poisson ou binomiale négative.

\subsection{Modèle hurdle}

\begin{align*}
P(Y=y) = \begin{cases}
1-p(x) & y=0\\
p(x)\,\dfrac{e^{-\lambda(z)}\lambda(z)^y}{y!\,(1-e^{-\lambda(z)})} & y=1,2,\ldots
\end{cases}
\end{align*}
On peut utiliser un modèle logistique pour $p(x)$~: $\text{logit}(p(x)) = x^T\beta$ et un modèle log-linéaire pour $\lambda(z)$~: $\log(\lambda(z)) = z^T\gamma$.

\textbf{Moments}~:
$$E(Y_i) = p(x_i)\,\frac{\lambda_i}{1-e^{-\lambda_i}}, \qquad Var(Y_i) = p(x_i)\,\frac{\lambda_i + \lambda_i^2}{1-e^{-\lambda_i}} - (E(Y_i))^2.$$

\textbf{Propriété clé}~: la log-vraisemblance se sépare en $l(\beta) + l(\gamma)$ $\Rightarrow$ les deux composantes peuvent être ajustées séparément (paramètres orthogonaux).

\subsection{Modèle ZIP (zero-inflated Poisson)}

\begin{align*}
P(Y=y) = \begin{cases}
\pi(x) + (1-\pi(x))\,e^{-\lambda(z)} & y=0\\
(1-\pi(x))\,\dfrac{e^{-\lambda(z)}\lambda(z)^y}{y!} & y=1,2,\ldots
\end{cases}
\end{align*}

\textbf{Moments}~:
$$E(Y_i) = (1-\pi(x_i))\,\lambda_i, \qquad Var(Y_i) = (1-\pi(x_i))(\lambda_i + \pi(x_i)\lambda_i^2) > E(Y_i) \text{ si } \pi > 0.$$

La vraisemblance doit être maximisée conjointement ou par l'algorithme EM.

\subsection{Comparaison hurdle vs ZIP}

\begin{center}
\begin{tabular}{l|l|l}
& \textbf{Hurdle} & \textbf{ZIP} \\
\hline
Paramétrisation & orthogonale & non orthogonale \\
Ajustement & séparé & simultané (ou EM) \\
Interprétation & $p_i$ = prob.\ de présence & $\pi$ = prob.\ de mélange \\
\end{tabular}
\end{center}

\textbf{Validation}~: tests du rapport de vraisemblance, AIC. Pour le modèle hurdle, le cadre des MLG s'applique pleinement.

\newpage
%% ============================================================
%% CHAPITRE 4 : GEE
%% ============================================================
\section{Équations d'estimation généralisées (GEE)}

\subsection{Données longitudinales}

Dans les \textbf{études longitudinales}, les individus sont mesurés de façon répétée au fil du temps. L'indépendance est supposée entre les individus, mais pas entre les mesures d'un même individu.

\subsection{Approches possibles}

\begin{itemize}
\item \textbf{Modèles marginaux}~: $g(E(Y_{it})) = x_{it}^T \beta$ et $Var(Y_i) = V_i(\phi,\alpha)$.
\item \textbf{Modèles à effets aléatoires}~: $g(E(Y_{it}|\gamma_i)) = x_{it}^T \beta + z_{it}^T \gamma_i$.
\end{itemize}

\subsection{Modèles marginaux et GEE}

L'inférence est de type \textbf{moyenne populationnelle}. Les paramètres $\beta$ sont estimés en résolvant~:
$$\sum_{i=1}^n D_i^T V_i(\alpha)^{-1}(Y_i - \mu_i) = 0,$$
avec $D_i = \partial \mu_i/\partial \beta$, $V_i(\alpha) = \phi\, A_i^{1/2} R(\alpha) A_i^{1/2}$, et $R(\alpha)$ la \textbf{matrice de corrélation de travail}.

\subsection{Choix de la matrice de corrélation de travail $R(\alpha)$}

\begin{description}
\item[Indépendance~:] $R(\alpha) = I$ (revient à un MLG).
\item[Échangeable~:] $R(\alpha)_{t,t'} = \alpha$ pour tout $t \neq t'$.
\item[Autorégressive (AR)~:] $R(\alpha)_{t,t'} = \alpha^{|t-t'|}$.
\item[$m$-dépendance~:] corrélation jusqu'à distance temporelle $m$.
\item[Non structurée~:] $R(\alpha)$ totalement libre.
\end{description}

\subsection{Estimation de $\phi$ et $\alpha$}

On utilise une procédure itérative~: scoring de Fisher modifié pour $\beta$ et estimation par les moments de $\alpha$ et $\phi$. Par exemple~:
$$\hat{\phi} = \frac{\sum_i \sum_t \hat{r}_{it}^2}{N - (p+1)}, \quad \text{où } r_{it} = \frac{y_{it} - \mu_{it}}{\sqrt{v(\mu_{it})}}.$$

\subsection{Distribution asymptotique}

$\sqrt{n}(\hat{\beta} - \beta)$ est asymptotiquement normal avec variance $\Omega = \lim_{n\to\infty} n\,M^{-1}QM^{-1}$, où~:
$$M = \sum_i D_i^T V_i^{-1} D_i, \qquad Q = \sum_i D_i^T V_i^{-1} Var(Y_i) V_i^{-1} D_i.$$

\begin{remarque}
La convergence de $\hat{\beta}$ ne dépend que de la spécification correcte de la moyenne $\mu_i$, pas de la corrélation. Choisir $R$ proche de la vraie corrélation augmente l'efficacité.
\end{remarque}

L'estimateur de la variance (dit ``sandwich'') est $\hat{\Omega} = \hat{M}^{-1}\hat{Q}\hat{M}^{-1}$.

\subsection{Diagnostics et inférence}

\textbf{Résidus de Pearson}~: $\hat{r}_{it} = (y_{it} - \hat{\mu}_{it})/\sqrt{v(\hat{\mu}_{it})}$.

L'inférence est limitée par l'absence de fonction de vraisemblance (pas de test du rapport de vraisemblance, pas d'AIC standard).

\newpage
%% ============================================================
%% CHAPITRE 5 : MLGM (GLMM)
%% ============================================================
\section{Modèles linéaires généralisés mixtes (MLGM)}

\subsection{Motivation}

Avec des données longitudinales ou groupées, il faut tenir compte de la corrélation intra-grappe. Le MLGM étend le MLG en permettant aux coefficients de varier selon la grappe.

\subsection{Définition}

\textbf{Ingrédients}~:
\begin{itemize}
\item Réponse~: $Y_i = (Y_{i1},\ldots,Y_{in_i})$.
\item Matrice de plan $X_i$ pour les \textbf{effets fixes}, $Z_i$ pour les \textbf{effets aléatoires}.
\item $\gamma_i$ : effets aléatoires latents (non observés).
\end{itemize}

\textbf{Hypothèses}~:
\begin{itemize}
\item $Y_{it} \mid \gamma_i$ indépendants, de la famille exponentielle.
\item $\gamma_i \stackrel{iid}{\sim} \mathcal{N}(0, \Psi)$.
\end{itemize}

\textbf{Modèle conditionnel}~:
$$g(E(Y_{it} \mid \gamma_i)) = g(\mu_{it}) = x_{it}^T \beta + z_{it}^T \gamma_i, \qquad Var(Y_{it} \mid \gamma_i) = v(\mu_{it}).$$

\subsection{Cas gaussien avec lien identité}

$$Y_i = X_i\beta + Z_i\gamma_i + \epsilon_i, \quad \epsilon_i \sim \mathcal{N}(0,\sigma^2_\epsilon I), \quad \gamma_i \sim \mathcal{N}(0,\Psi).$$
Cela implique $Y_i \sim \mathcal{N}(X_i\beta,\, Z_i\Psi Z_i^T + \sigma^2_\epsilon I)$.

\subsection{Estimation}

On s'appuie sur la \textbf{vraisemblance marginale}~:
$$L(\beta,\phi,\Psi) = \prod_{i=1}^n \int \prod_{t=1}^{n_i} f(y_{it} \mid \gamma_i,\beta,\phi)\, f(\gamma_i \mid \Psi)\, d\gamma_i.$$

Cette intégrale ne peut en général pas être calculée en forme fermée. Approximations~:
\begin{itemize}
\item Quadrature de Gauss-Hermite
\item Approximation de Laplace (défaut dans \texttt{lme4})
\item Quasi-vraisemblance pénalisée (PQL) / marginale (MQL) -- peuvent produire des estimateurs biaisés
\item MCMC
\end{itemize}

\subsection{Interprétation des effets fixes}

Un MLGM est une analyse \textbf{spécifique au sujet}~: $\mu_{it} = g^{-1}(x_{it}^T\beta + z_{it}^T\gamma_i)$.

La moyenne populationnelle est~:
$$E(Y_{it}) = \int g^{-1}(x_{it}^T\beta + z_{it}^T\gamma_i)\,dF(\gamma_i) \neq g^{-1}(x_{it}^T\beta),$$
sauf dans le cas gaussien avec lien identité.

\subsection{Prédiction des effets aléatoires}

Les logiciels fournissent les \textbf{modes conditionnels}~: valeurs maximisant $f(\gamma_i \mid y)$. Les meilleurs prédicteurs seraient $\hat{\gamma}_i = E(\gamma_i \mid y_i)$.

\textbf{Prédiction}~:
\begin{itemize}
\item Nouveau point dans un groupe existant~: $g^{-1}(x_{it}^T\hat{\beta} + z_{it}^T\hat{\gamma}_i)$.
\item Nouveau groupe~: $g^{-1}(x_{it}^T\hat{\beta})$.
\end{itemize}

\newpage
%% ============================================================
%% CHAPITRE 6 : Méthodes robustes
%% ============================================================
\section{Méthodes robustes}

\subsection{Motivation}

Les analyses statistiques classiques peuvent être fortement affectées par les valeurs aberrantes. La \textbf{statistique robuste} considère que~:
\begin{itemize}
\item les modèles sont au mieux des approximations,
\item des déviations par rapport aux hypothèses sont presque toujours présentes.
\end{itemize}

On suppose que la distribution provient d'un \textbf{voisinage du modèle postulé}~:
$$F_\epsilon = (1-\epsilon)F_\theta + \epsilon G,$$
où $G$ est une distribution arbitraire de contamination.

\subsection{Outils de mesure de la robustesse}

\subsubsection{Fonction d'influence}
$$IF(z; T, F) = \lim_{\epsilon \to 0} \frac{T(F_\epsilon) - T(F)}{\epsilon},$$
où $F_\epsilon = (1-\epsilon)F + \epsilon\Delta_z$. Mesure l'effet d'une contamination infinitésimale au point $z$. Une \textbf{fonction d'influence bornée} est souhaitable.

\subsubsection{Point de rupture (\textit{breakdown point})}
Pourcentage maximal de la minorité tel que celle-ci n'a qu'une influence limitée. Plus le point de rupture est \textbf{grand}, plus l'estimateur est robuste.

\subsection{M-estimation}

Les \textbf{estimateurs M} sont définis par~: $\sum_{i=1}^n \Psi(y_i, \theta) = 0$.

Propriétés~: $\sqrt{n}(\hat{\theta}-\theta) \sim \mathcal{N}(0, M^{-1}QM^{-1})$ et $IF \propto \Psi$.

L'EMV pour les MLG est un estimateur M avec $\Psi$ non bornée $\Rightarrow$ \textbf{non robuste}.

\subsection{MLG robustes}

On introduit~:
$$\sum_{i=1}^n \left[\psi(r_i)\,w(x_i)\,\frac{\mu_i'}{v^{1/2}(\mu_i)} - a(\beta)\right] = 0,$$
où $\psi$ contrôle les grandes déviations en $y$, $w(x_i)$ réduit l'influence des points à levier, et $a(\beta)$ assure la consistance de Fisher.

\textbf{Fonction de Huber}~: $\psi_c(r) = r\cdot\min(1, c/|r|)$, où $c$ contrôle le compromis robustesse/efficacité (en pratique $c \in [1,2]$).

\textbf{Poids $w(x_i)$}~: basés sur les éléments diagonaux de la matrice chapeau ou les distances de Mahalanobis robustes.

\subsection{Sélection de variables robuste}

On utilise la \textbf{quasi-vraisemblance robuste} et la statistique $\Lambda_{QM}$ (différence de quasi-vraisemblances). Sa distribution sous $H_0$ est une combinaison linéaire de $\chi^2$.

\subsection{GEE robustes}

Version robuste des GEE~:
$$\sum_{i=1}^n D_i^T \Gamma_i^T V_i^{-1}(\psi_i - c_i) = 0,$$
où $\psi_i = W_i(Y_i - \mu_i)$ avec $W_i$ diagonale pour limiter l'influence.

\subsection{Modèle hurdle robuste}

Combine une régression binaire robuste pour séparer zéros/non-zéros et un MLG robuste pour la Poisson tronquée.

\newpage
%% ============================================================
%% CHAPITRE 7 : MAG (GAM)
%% ============================================================
\section{Modèles additifs généralisés (MAG)}

\subsection{Idée générale}

On remplace la partie linéaire du MLG par des fonctions non paramétriques~:
$$g(\mu_i) = \alpha + f_1(x_{i1}) + \cdots + f_p(x_{ip}).$$

\begin{center}
\textit{Laissons les données nous montrer la forme fonctionnelle appropriée.}
\end{center}

\subsection{Lissage univarié}

\subsubsection{Fonctions de base}
On représente $f(x)$ par~: $f(x) = \sum_{k=1}^K b_k(x)\beta_k$. Le modèle devient $Y = X\beta + \epsilon$, ajusté par moindres carrés.

Exemples de bases~: polynomiale, linéaire par morceaux (tente), splines.

\subsubsection{Lissage pénalisé}

Plutôt que sélectionner $K$, on choisit $K$ ``suffisamment grand'' et on pénalise l'oscillation~:
$$\hat{\beta} = \arg\min_\beta \left\{\|y - X\beta\|^2 + \lambda\,\beta^T S\beta\right\},$$
avec solution explicite~:
$$\hat{\beta} = (X^TX + \lambda S)^{-1}X^Ty.$$

Le \textbf{paramètre de lissage} $\lambda$ contrôle le compromis régularité/fidélité~:
\begin{itemize}
\item $\lambda \to \infty$~: estimation linéaire de $f$.
\item $\lambda = 0$~: régression non pénalisée (risque de surajustement).
\end{itemize}

\subsubsection{Matrice d'influence et degrés de liberté effectifs}

$$\hat{y} = Ay, \quad A = X(X^TX + \lambda S)^{-1}X^T.$$

Les \textbf{degrés de liberté effectifs} (DLE) sont~:
$$\text{DLE} = Tr(F), \quad F = (X^TX + \lambda S)^{-1}X^TX.$$

La pénalisation induit un \textbf{biais}~: $E(\hat{\beta}) = F\beta \neq \beta$.

\subsubsection{Interprétation bayésienne}

La pénalisation est équivalente à une loi a priori gaussienne~:
$$\beta \sim \mathcal{N}(0, S^{-}\sigma^2/\lambda).$$

On obtient la loi a posteriori~:
$$\beta \mid y \sim \mathcal{N}(\hat{\beta},\, \sigma^2(X^TX + \lambda S)^{-1}).$$

\subsubsection{Choix de $\lambda$ par validation croisée}

\textbf{Validation croisée ordinaire (OCV)}~:
$$OCV(\lambda) = \frac{1}{n}\sum_{i=1}^n\left(\frac{y_i - \hat{f}_\lambda(x_i)}{1-A_{ii}}\right)^2.$$

\textbf{Validation croisée généralisée (GCV)}~:
$$GCV(\lambda) = \frac{n\sum_i(y_i - \hat{f}_\lambda(x_i))^2}{(n - Tr(A))^2}.$$

\subsection{Modèles additifs (plusieurs prédicteurs)}

$$y = \alpha + f_1(x_1) + f_2(x_2) + \cdots + f_p(x_p) + \epsilon,$$
avec la contrainte d'identifiabilité $\sum_i f_j(x_{ji}) = 0$ pour tout $j$.

On résout~:
$$\hat{\beta} = \arg\min_\beta\left\{\|y - X\beta\|^2 + \sum_{j=1}^p \lambda_j\beta^TS_j\beta\right\},$$
avec solution~: $\hat{\beta} = (X^TX + \sum_j \lambda_j S_j)^{-1}X^Ty$.

\textbf{Distribution a posteriori}~: $\beta \mid y \sim \mathcal{N}(\hat{\beta}, V_\beta)$ avec $V_\beta = \sigma^2(X^TX + \sum_j \lambda_j S_j)^{-1}$.

\subsubsection{Tests d'hypothèses}

Pour tester $H_0: \beta_j = 0$ (sous-vecteur de taille $p_j$)~:
$$\hat{\beta}_j^T V_{\beta_j}^{-1}\hat{\beta}_j / p_j \sim F_{p_j, n-p} \quad (\phi \text{ inconnu}), \qquad \hat{\beta}_j^T V_{\beta_j}^{-1}\hat{\beta}_j \sim \chi^2_{p_j} \quad (\phi \text{ connu}).$$

\subsubsection{Intervalles de confiance pour $f_j$}
$$\hat{f}_{ji} \pm z_{\alpha/2}\sqrt{v_i}, \quad v = \text{diag}(\tilde{X}V_\beta\tilde{X}^T).$$
Les intervalles de crédibilité bayésiens possèdent de bonnes propriétés de couverture fréquentistes.

\subsection{Modèles additifs généralisés (GAM)}

Comme pour les MLG, on combine la structure additive avec les distributions de la famille exponentielle~:
$$g(\mu_i) = \alpha + \sum_{j=1}^p f_j(x_{ij}) = \alpha + \sum_{j=1}^p X^{(j)}\delta^{(j)}.$$

On minimise~:
$$\hat{\beta} = \arg\min_\beta\left\{D(y,\mu) + \sum_{j=1}^p \lambda_j \beta^T S_j \beta\right\},$$
où $D(y,\mu)$ est la déviance.

\subsubsection{Algorithme de score local}

À l'étape $(m-1)$~:
\begin{enumerate}
\item Construire la variable dépendante ajustée $z_i^{(m-1)} = \eta_i^{(m-1)} + (y_i - \mu_i^{(m-1)})\left(\frac{\partial \eta_i}{\partial \mu_i}\right)^{(m-1)}$.
\item Construire les poids $w_i^{(m-1)} = \frac{1}{Var(Y_i)^{(m-1)}}\left(\frac{\partial \mu_i}{\partial \eta_i}\right)^2_{(m-1)}$.
\item Ajuster un modèle additif pondéré à $z^{(m-1)}$.
\item Répéter jusqu'à convergence.
\end{enumerate}

\subsubsection{Choix de $\lambda_1,\ldots,\lambda_p$}

\begin{itemize}
\item Si $\phi$ connu~: \textbf{UBRE} $= D(y,\hat{\mu}) + 2\phi\,Tr(F)/n$ (aussi appelé $C_p$ de Mallows).
\item Si $\phi$ inconnu~: \textbf{GCV} $= \frac{n\,D(y,\hat{\mu})}{(n-Tr(F))^2}$.
\end{itemize}

\subsubsection{Comparaison de modèles}

\begin{itemize}
\item \textbf{AIC}~: $-2l(\hat{\beta};y) + 2\phi\,Tr(F)$. Plus petit = meilleur.
\item \textbf{GCV}~: estime l'erreur de prédiction. Plus petit = meilleur.
\item Proportion de déviance expliquée.
\end{itemize}

\newpage
\printbibliography

\end{document}
